Naše intuice o celých číslech je obvykle řízena představou, že to jsou čísla,
která navíc k~obvyklým počtům předmětů (vyjádřených
\hyperref[chap:prirozena-cisla]{čísly přirozenými}) umožňují uchopit představu
\emph{ztráty}, \emph{dluhu} či \emph{nedostatku} pomocí čísel záporných.

Tato intuice není zcestná, ale rovněž není \emph{úplná}. Totiž, fakt, že záhadné
znaménko $-$ před číslem způsobí, že místo přebytku vyjadřuje nedostatek, je
snad přirozená, ale jest, tvrdíme, méně přirozené, že by třeba nedostatek
nedostatku měl znamenat přebytek. To, že mi nic nechybí, nezbytně neznamená, že
něco mám, že? \emph{Násobení} celých čísel je ještě méně selsky uchopitelné.
Jaký dává smysl, že součinem dvou záporů je klad?

My zde vyrobíme celá čísla poněkud jinak, než byste snad, čtenáři, čekali.
Nikoli jako čísla záporná a kladná, anobrž jako čísla vyjadřující
\textbf{vzdálenosti mezi přirozenými čísly}. K~tomu vysvětlíme, jak přesně
míníme slovo \uv{vzdálenost}. Například, vzdálenost mezi čísly $2$ a $5$ je
zřejmě $3$, ovšem stejná vzdálenost je i mezi čísly $5$ a $2$. S takovýmto
popisem vzdálenosti nejsme spokojeni, neb v sobě nenese informaci o tom, které
číslo je menší. Rozšíříme jej proto ještě o \emph{směr}. Když bude první z čísel
menší, dáme tento fakt najevo umístěním symbolu $-$ před jejich vzdálenost.
\textbf{Pozor!} V našem pojetí neznačí toto znaménko žádný \uv{zápor}, leč pouze
skutečnost, že směr od $2$ k $5$ je opačný od směru od $5$ k $2$. Znovu:
řekneme, že vzdálenost (či \emph{rozdíl}) mezi $2$ a $5$ je $-3$, zatímco
vzdálenost mezi $5$ a $2$ je $3$; symbol $-$ nese význam \emph{směru}, nikoli
záporu.

Než se vrhneme do formalisace právě vyřčených konceptů vzdálenosti a směru,
potřebujeme rozšířit své obzory o tzv. \emph{třídy ekvivalence}. Co jsou zač a k
čemu jich použijeme se dozvíte záhy.

\section{Třídy ekvivalence}
\label{sec:tridy-ekvivalence}

V \myref{pododdíle}{sssec:relace} jsme definovali
\hyperref[def:relace-ekvivalence]{relaci ekvivalence}. Je to relace, která v
podstatě umožňuje považovat dva prvky množiny za \emph{stejné}, za předpokladu,
že sdílejí jisté vlastnosti. Intuitivně využíváme ekvivalence neustále, téměř
kdykoli pojmenováváme objekty. Žádné dva objekty, které vnímáme, nejsou přímo
\emph{stejné}, ale přiříkáme jim tatáž jména, protože jsou \uv{dostatečně
podobné}, abychom je mohli označit stejně. Kupříkladu, starou dřevěnou vysokou
židli i novou malou kovovou židli jsme právě nazvali \emph{židlí}, neboť obě
mají dostatek společných charakteristik a plní stejný účel, přestože se jistě
nejedná o týž předmět.

Příklad židle s sebou nese další samozřejmou, leč užitečnou, myšlenku: svůj svět
přirozeně rozdělujeme na skupiny předmětů, které považujeme za \uv{stejné}. V
matematice se taková \emph{skupina podobných předmětů} nazývá \alert{třída
ekvivalence}. Je to množina všech prvků, které jsou vzájemně ekvivalentní.

\begin{definition}[Třída ekvivalence]
  \label{def:trida-ekvivalence}
  Ať $e$ je ekvivalence na množině $a$. Pro libovolný $x \in a$ říkáme všem
  prvkům z $a$, které jsou ekvivalentní s $x$, \emph{třída ekvivalence příslušná
  $x$} a značíme ji $[x]_e$. Formálně, pro každé $x \in a$ je
  \[
    [x]_e = \{y \in a \mid (x, y) \in e\}.
  \]
  Prvku $x$ pak říkáme \emph{representant} této třídy ekvivalence.
\end{definition}
\begin{example}
  \label{exam:trida-ekvivalence-1}
  Ať $l$ je množina lidí a $e$ je ekvivalence \uv{být stejně starý} (je to
  opravdu ekvivalence?). Když je $\text{Adam} \in l$ člověk, kterému je 29 let.
\end{example}

